% STATUS: DRAFT
\chapter{Local algebras are type III$_1$}\label{ch:typeIII_local}

\begin{physmotiv}
We now collect the evidence for the claim on which the whole book rests: algebras of local regions in QFT are of type III$_1$. This explains in one stroke the divergence of entanglement entropy, the absence of local pure states, and the Reeh--Schlieder property, and tells us exactly what must change when gravity is included: we need a trace.
\end{physmotiv}

\section{The statement}

\begin{center}
\emph{In a (reasonable) local QFT, the algebra $\A(O)$ of a bounded region with non-trivial causal complement is the hyperfinite factor of type III$_1$.}
\end{center}

Parts of this are theorems; other parts require assumptions. We list the logic, separating what is proven under general axioms from what needs more.

\begin{enumerate}
\item \textbf{Cyclic and separating vacuum} (Reeh--Schlieder, \cref{ch:rs_bw}) $\Rightarrow$ modular theory applies, and $\omega_0|_{\A(O)}$ is a faithful normal state.
\item \textbf{No minimal projections, no trace.} A type I factor has pure normal states. Pure normal states are incompatible with analyticity in the translations (the energy-momentum spectrum condition), so $\A(O)$ is not type I. Under conditions that include the \emph{split property} and scaling-type assumptions one shows that no finite trace exists; local algebras are of type III (Fredenhagen, Buchholz, and others; see the literature).
\item \textbf{Modular spectrum.} For wedges, $\Delta^{it}$ is the boost and $\spec\Delta=[0,\infty)$ because the boost generator has spectrum $\R$ (\cref{ch:connes_class}, exercise). Together with factoriality this shows the wedge algebra is type III$_1$ (Driessler; Longo; Fredenhagen; see Buchholz--D'Antoni--Longo for the passage to double cones).
\item \textbf{Uniqueness.} If the net satisfies nuclearity/split conditions and has a non-trivial scaling limit, $\A(O)$ is the unique \emph{hyperfinite} type III$_1$ factor $R_\infty$ (Haagerup's theorem identifies all injective type III$_1$ factors; Buchholz--D'Antoni--Longo establish injectivity under their hypotheses).
\end{enumerate}

The upshot is that, for the theories we care about (free fields, conformal theories, and any theory with a good scaling limit in the UV), local algebras are \emph{all isomorphic} to one single abstract algebra: $R_\infty$, the Araki--Woods factor of \cref{ch:connes_class}. Differences between theories are not in the isomorphism class of $\A(O)$ but in how the family of algebras sits in $\Bnd{\HH}$.

\section{Consequences}

\begin{center}
\renewcommand{\arraystretch}{1.3}
\begin{tabular}{@{}p{4.6cm}p{8cm}@{}}
\toprule
Property of type III$_1$ & Physics\\
\midrule
No trace; no density matrix & Region entropy $-\Tr\rho\log\rho$ ill defined; UV divergence $\sim\mathrm{Area}/\epsilon^{d-2}$\\
No pure normal states & No local state is pure; any state is entangled with the environment\\
All projections $e\sim\id$ & Local operators can generate any state from the vacuum (Reeh--Schlieder)\\
Modular spectrum $[0,\infty)$; flow outer & The boost acts as an outer automorphism: Unruh/Hawking temperature without a Hamiltonian in the region\\
Maximal Bell violation & Entanglement is as strong as quantum mechanics allows (\cref{ch:subsystems})\\
Entanglement dense & Product states do not exist as normal states of $\A(R)\vee\A(L)$\\
\bottomrule
\end{tabular}
\end{center}

\begin{whatwrong}
``Type III$_1$'' is an exact statement about the continuum algebra. The lattice, with finite-dimensional local Hilbert spaces, is type I, and the area-law entropy there is finite and cutoff dependent. The continuum limit loses the factorization, not just the finiteness. Equally, the absence of a trace does \emph{not} mean entropy differences are meaningless: relative entropy and the first law survive (\cref{ch:modham}).
\end{whatwrong}

\section{What gravity must supply}

The structure theory of \cref{ch:connes_class,ch:crossed,ch:takesaki} tells us the minimal modification that yields a well-defined entropy: take the crossed product by the modular flow, obtaining a type II$_\infty$ algebra with a trace. The remaining question, of \emph{physics}, is why gravity should do this. The answer (Part VII) is that the modular flow of a horizon is a symmetry whose generator is a \emph{constraint} in a diffeomorphism-invariant theory: gauge-invariant observables must commute with it, which means they must be dressed to a physical clock, and the algebra of dressed observables is the crossed product.

\begin{keyidea}
Local algebras of QFT are (in nice theories) the hyperfinite III$_1$ factor. Everything about vacuum entanglement that is divergent or non-factorizing is a symptom of this. The minimal cure is the crossed product, and gravity is expected to provide it.
\end{keyidea}

\section*{Exercises}
\begin{exercise}
Show that a type I$_\infty$ factor $\Bnd K\otimes\id\subset\Bnd{K\otimes K'}$ would give $\A(O)$ a pure normal state, and explain why this contradicts the fact that the vacuum is cyclic for $\A(O')=\A(O)'$.
\end{exercise}
\begin{exercise}
Starting from the free-field lattice entropy $S=\frac13\log(L/\epsilon)$ in $1+1$ dimensions, explain why the $\epsilon$-dependence cannot be removed by any state-independent subtraction that is covariant in the continuum, in the language of type III.
\end{exercise}
\begin{exercise}
List which of the following are type III$_1$: the wedge algebra of a free field; the algebra of a double cone in a CFT; the algebra of the whole Minkowski space in the vacuum representation; the algebra of an ordinary quantum mechanical particle. Justify each.
\end{exercise}
